\documentclass[
	12pt,
	a4paper,
	bibliography=totoc,
	numbers=endperiod,
	parskip=half,
]{scrartcl}

\usepackage[T1]{fontenc}
\usepackage[utf8]{inputenc}

\usepackage{libertinus}
\usepackage[varqu,scaled=0.95]{inconsolata}
\usepackage{newtxmath}

\usepackage[english]{babel}

% Use less space for the page number
\usepackage[margin=2.5cm,footskip=36pt]{geometry}
\usepackage{graphicx}
\usepackage[htt]{hyphenat}
\usepackage{listings}
\usepackage{microtype}
\usepackage{subcaption}
% Required for listings's upquote
\usepackage{textcomp}
\usepackage{upquote}
\usepackage{xcolor}

\usepackage[hyphens]{url}
\usepackage[hidelinks]{hyperref}
\usepackage[capitalise,noabbrev]{cleveref}

\usepackage[color=ovgu-orange]{todonotes}

\usepackage{lipsum}

\graphicspath{{./figures/}}

\definecolor{ovgu-blue}{HTML}{0068B4}
\definecolor{ovgu-darkgray}{HTML}{606060}
\definecolor{ovgu-lightgray}{HTML}{C0C0C0}
\definecolor{ovgu-orange}{HTML}{F39100}
\definecolor{ovgu-purple}{HTML}{7A003F}

\lstset{
	basicstyle=\ttfamily,
	commentstyle=\color{ovgu-darkgray},
	keywordstyle=\color{ovgu-blue},
	numberstyle=\ttfamily\color{ovgu-darkgray},
	stringstyle=\color{ovgu-purple},
	rulecolor=\color{ovgu-lightgray},
	frame=single,
	numbers=left,
	language=C,
	breaklines=true,
	breakatwhitespace=true,
	postbreak=\hbox{$\hookrightarrow$ },
	showlines=true,
	showstringspaces=false,
	upquote=true,
	tabsize=4,
	gobble=0,
	captionpos=b,
	abovecaptionskip=\medskipamount,
}

\titlehead{\centering\includegraphics[width=0.66\textwidth]{OVGU-INF}}

\title{Title}
\subtitle{Subtitle}

\author{
First Author\\
{\large\href{mailto:first.author@ovgu.de}{\nolinkurl{first.author@ovgu.de}}}
\and
Second Author\\
{\large\href{mailto:second.author@ovgu.de}{\nolinkurl{second.author@ovgu.de}}}
}

\date{\today}

\begin{document}

\maketitle

\begin{abstract}
\lipsum[1]

This report template is available at \url{https://github.com/parcio/templates} and consists of \cref{sec:introduction,sec:background,sec:evaluation,sec:conclusion}.
\end{abstract}

\newpage

\tableofcontents

\newpage

\section{Introduction}
\label{sec:introduction}

\begin{figure}[ht]
	\centering
	\begin{subfigure}{0.45\textwidth}
		\centering
		\includegraphics[width=0.9\textwidth]{OVGU-INF}
		\caption{Left}
		\label{fig:left}
	\end{subfigure}
	\begin{subfigure}{0.45\textwidth}
		\centering
		\includegraphics[width=0.9\textwidth]{OVGU-INF}
		\caption{Right}
		\label{fig:right}
	\end{subfigure}
	\caption{Caption}
	\label{fig:both}
\end{figure}

You can refer to the subfigures (\cref{fig:left,fig:right}) or the figure (\cref{fig:both}).

\section{Background}
\label{sec:background}

You can comfortably reference literature~\cite{DBLP:journals/superfri/DuweLMSF0B020}.\footnote{This is a footnote.}

\begin{table}[ht]
	\centering
	\begin{tabular}{|l|c|r|}
		\hline
		\textbf{Header 1} & \textbf{Header 2} & \textbf{Header 3} \\
		\hline
		\hline
		Row 1 & Row 1 & Row 1 \\
		Row 2 & Row 2 & Row 2 \\
		\hline
	\end{tabular}
	\caption{Caption}
	\label{tab:table}
\end{table}

You can also refer to tables (\cref{tab:table}).

\subsection{Math}

\[
	E = m c^2
\]

\section{Evaluation}
\label{sec:evaluation}

\begin{lstlisting}[caption=Caption,label=lst:hello-world]
printf("Hello world!\n");

// Comment
for (int i = 0; i < m; i++) {
	for (int j = 0; j < n; j++) {
		sum += 'a';
	}
}
\end{lstlisting}

You can also refer to listings (\cref{lst:hello-world}).

\section{Conclusion}
\label{sec:conclusion}

\todo[inline]{FIXME}

\lipsum[1-3]
\todo{FIXME: remove this}

\newpage

\bibliographystyle{apalike}
\bibliography{report}

\end{document}
